\title{Direct Preference Optimization: Your Language Model is Secretly a Reward Model}

\begin{abstract}
Large-scale unsupervised language models (LMs) acquire extensive general knowledge and some capabilities in reasoning, but their behavior is challenging to direct due to the nature of their purely unsupervised training processes. Contemporary approaches to making these models more controllable typically involve collecting human judgments that rank model outputs and adapting the pretrained LM to reflect these preferences, most often utilizing reinforcement learning from human feedback (RLHF). Nevertheless, RLHF introduces considerable complexity and instability: it involves first constructing a reward model to encode human preferences, and subsequently optimizing the pretrained LM with reinforcement learning so as to maximize this learned reward while preventing excessive deviation from the initial model. This work proposes a novel reparameterization of the RLHF reward model, making it possible to directly derive the optimal policy in closed form. This enables us to address the RLHF objective using only a straightforward classification loss. We term this approach \textit{Direct Preference Optimization} (DPO); it is both robust and efficient, requiring no sampling from the LM during the fine-tuning process or extensive hyperparameter tuning. Experimental results indicate that DPO enables effective alignment of LMs with human preferences, performing on par with or better than current techniques. Specifically, DPO surpasses RLHF based on PPO for sentiment control in model generations, and provides comparable or enhanced performance in tasks such as summarization and single-turn dialogue, while significantly reducing implementation and training complexity.
\end{abstract}

\section{Introduction}
Large unsupervised language models (LMs), which are trained on massive corpora, often demonstrate unexpected and impressive capabilities~\citep{chowdhery2022palm, brown2020language, touvron2023llama,bubeck2023sparks}. Nevertheless, because these models are trained on text authored by humans whose goals, motivations, and skill levels vary greatly, the models inevitably absorb a heterogeneous set of behaviors and knowledge. Some of these inherited traits are not always desirable for direct replication. For example, an AI coding assistant should be able to \textit{recognize} prevalent programming errors to effectively correct them, but, ideally, its code generation should favor the rarer, high-caliber coding practices found in its training distribution. In a similar vein, even though a language model should be \textit{informed} about a widely held misconception—such as one believed by 50\% of people—it is inappropriate for the model to reproduce this false belief in 50\% of its responses. Put simply, it is essential to carefully control the model’s \emph{preferred behaviors and answers} in order to ensure that deployed AI systems remain safe, effective, and aligned with human values~\citep{ouyang2022training}. Although reinforcement learning (RL) is often used to align LMs with human preference signals, we demonstrate that the optimization objective used in existing RL-based methods can, in fact, be achieved exactly via a straightforward binary cross-entropy objective, thus offering a much more streamlined approach to preference learning.

In general, present-day techniques impart targeted behaviors to language models by leveraging curated collections of human preference data that reflect safety and helpfulness criteria. This phase of preference learning follows the initial, large-scale unsupervised pre-training on vast text corpora. While the simplest way to conduct preference learning involves supervised fine-tuning on exemplary human responses, reinforcement learning from human or AI feedback (RLHF/RLAIF; \citep{christiano2017deep,bai2022constitutional}) has emerged as the most successful methodology. RLHF methods begin by training a reward model on datasets of ranked human preferences, subsequently employing RL to adjust the language model’s output policy so that it maximizes the expected reward—while constraining it from deviating too far from the pre-trained policy. Despite RLHF’s success in producing language models with advanced dialogue and coding performance, the RLHF process introduces substantial complexity over standard supervised approaches: it requires the training and coordination of multiple models and the recurrent sampling from LM policies during optimization, leading to high computational demands.

This work introduces a method for optimizing language models to conform with human preferences, bypassing both explicit reward modeling and reinforcement learning. We present \textit{{Direct Preference Optimization} (DPO)}, an approach that achieves the same optimization goal as traditional RLHF techniques—maximizing reward under a KL-divergence regularization—while offering a more accessible implementation and training process. Conceptually, DPO operates by elevating the relative log-likelihood of favored responses compared to those not preferred, employing a sample-specific and adaptive importance weight to circumvent model collapse issues commonly seen with straightforward probability ratio approaches. Similar to prior algorithms, DPO leverages a theoretical preference framework (e.g., the Bradley-Terry model; \cite{bradley1952rankanalysis}) to assess how effectively a reward function captures observed preference data. Distinct from conventional strategies, which first train a reward model using a preference loss and then train a policy to optimize that learned reward, DPO uses a change of variables to express the preference loss directly as a function of the policy itself. Utilizing human preference datasets for model responses, DPO directly optimizes the policy by minimizing a binary cross-entropy objective, thereby producing the policy that implicitly aligns with the reward best fitting the empirical preferences.

The central innovation of this study is Direct Preference Optimization (DPO), which eliminates the need for reinforcement learning in the preference-driven training of language models. Empirical results indicate that DPO performs on par with, or surpasses, prevailing techniques such as PPO-based RLHF, across various preference-sensitive tasks like sentiment control, summarization, and dialogue, and scales successfully to language models with up to 6B parameters.

\section{Related Work}

Recent advances in self-supervised language models, particularly as their scale increases, have demonstrated their capability to address certain tasks with zero-shot \citep{radford2019language} or few-shot prompting approaches \citep{gpt3,megatron,chowdhery2022palm}. Nonetheless, their effectiveness on downstream applications and their alignment with user expectations are greatly enhanced by fine-tuning on curated datasets that pair instructions with human-generated completions \citep{mishra-etal-2022-cross,sanh2022multitask,chung2022scaling,thoppilan2022lamda}. This process, often called `instruction-tuning,' not only enables LLMs to better generalize to novel instructions that were absent from the fine-tuning corpus but also typically enhances their overall utility \citep{chung2022scaling}. While instruction-tuning has demonstrated substantial success, collecting \textit{relative} assessments of completion quality from humans is often simpler than assembling datasets of expert-crafted demonstrations. As a result, later studies have explored the fine-tuning of LLMs using collections of human preference data, leading to improved outcomes in domains such as translation \citep{kreutzer-etal-2018-reliability}, summarization \citep{stiennon2022learning,ziegler2020finetuning}, story generation \citep{ziegler2020finetuning}, and instruction following \citep{ouyang2022training,ramamurthy2023is}. These approaches typically begin by learning a neural reward model aligned with human preference datasets using models such as Bradley-Terry \citep{bradley1952rankanalysis}. The language model is subsequently fine-tuned to optimize this reward function via reinforcement learning algorithms, including REINFORCE \citep{williams1992reinforce}, proximal policy optimization (PPO; \cite{schulman2017proximal}), or related adaptations \citep{ramamurthy2023is}. An adjacent line of research utilizes LLMs, which have already been tuned to follow instructions with human oversight, to generate new synthetic preference data targeting qualities like safety and harmlessness \citep{bai2022constitutional}. This synthetic data generation uses minimal direct human input, relying instead on textual rubrics provided to the LLM for annotation. Collectively, these strategies represent an intersection of two research directions: one focusing on training language models through reinforcement learning for diverse targets~\citep{Ranzato2015SequenceLT,paulus2018a,wu2018learning}, and the other on algorithms designed to learn from human preferences \citep{christiano2017deep,kupcsik2018learning}. Although leveraging comparative human preferences offers several advantages, fine-tuning large-scale language models with reinforcement learning introduces significant practical difficulties; this work offers a principled method for optimizing with relative preferences that circumvents the need for reinforcement learning.

Beyond language-focused applications, the study of learning policies from preferences extends into both bandit and reinforcement learning domains, giving rise to various methodologies. One such paradigm is the contextual dueling bandit (CDB; \cite{yue2012karmed,dudik2015contextual}), in which policy optimization relies on comparative preferences or rankings of actions rather than direct reward signals. Since absolute rewards are unavailable in this setting, the theoretical underpinning of CDBs involves replacing the notion of an optimal policy with the \textit{von Neumann winner}—defined as a policy whose probability of outperforming any other policy is at least 50\% on expectation \citep{dudik2015contextual}. Importantly, while CDBs presume the availability of preference feedback in an online fashion, real-world human preference learning often deals with a static, pre-collected dataset comprising preference-labeled action pairs \citep{yan2022human}. In a related vein, \textit{preference-based RL} (PbRL) explores learning from binary preferences produced by an \textit{unknown} `scoring' or reward function instead of explicit rewards \citep{BusaFekete2014,ruiz2023dueling}. A range of PbRL algorithms exist—including off-policy methods that reuse stored preference data—but most entail a two-stage process: first modeling the latent scoring (reward) function, followed by policy optimization with respect to this estimated function \citep{jain2013learning,BusaFekete2014,christiano2017deep,sadigh2017active,kupcsik2018learning}. By contrast, our approach employs a unified, direct optimization procedure that trains policies to align with preference data in a single stage.

\section{Preliminaries}\label{section:prelims}

We briefly review the RLHF process described in \citeauthor{ziegler2020finetuning} (and subsequent works such as \citep{stiennon2022learning, bai2022training, ouyang2022training}). This process is conventionally divided into three main stages: (1) supervised fine-tuning (SFT); (2) acquisition of preference data and reward model training; and (3) reinforcement learning-based optimization.

\textbf{SFT}: The RLHF pipeline commences with supervised fine-tuning, wherein a pre-trained language model undergoes further supervised adaptation using task-specific, high-quality annotated data (e.g., dialogue systems, summarization tasks), resulting in a fine-tuned policy denoted as $\pi^\text{SFT}$.

\textbf{Reward Modelling Phase}: In the subsequent phase, the SFT model is conditioned on input prompts $x$ to produce two outputs $(y_1, y_2)\sim \pi^\text{SFT}(y \mid x)$. Human annotators are then asked to select a preferred response, yielding a label of the form $y_w\succ y_l \mid x$, with $y_w$ and $y_l$ identifying the chosen and rejected completions, respectively. The underlying mechanism generating these preferences is modeled as a latent reward function $r^*(y, x)$, which is assumed unknown. To capture the human preference structure, a variety of models can be used; among these, the Bradley-Terry (BT) model \cite{bradley1952rankanalysis} is commonly adopted, although extensions like the Plackett-Luce ranking model \citep{plackett1975analysis, luce2012individual} can accommodate situations where multiple ranked options are provided. The BT model expresses the human preference distribution $p^*$ as follows:

\begin{equation}\label{eq:bradley-terry}
    p^*(y_1\succ y_2 \mid x)=\frac{\exp\left(r^*(x, y_1)\right)}{\exp\left(r^*(x, y_1)\right) + \exp\left(r^*(x, y_2)\right)}.
\end{equation}

Given a fixed dataset of comparisons $\mathcal{D}=\bigl\{x^{(i)}, y_w^{(i)}, y_l^{(i)}\bigr\}_{i=1}^N$ sampled according to $p^*$, we may define a parameterized reward function $r_{\phi}(x, y)$ and infer its parameters by maximizing the likelihood. Casting this setup as a binary classification task leads to the negative log-likelihood objective:

\begin{equation}\label{eq:reward_model}
    \mathcal{L}_R(r_{\phi}, \mathcal{D}) = -\mathbb{E}_{(x, y_w, y_l)\sim \mathcal{D}}\bigl[\log \sigma(r_{\phi}(x, y_w)- r_{\phi}(x, y_l))\bigr]
\end{equation}

where $\sigma$ denotes the sigmoid (logistic) function. For language models, the network $r_{\phi}(x, y)$ is commonly initialized from the SFT policy $\pi^\text{SFT}(y \mid x)$, and a linear projection is appended to the output of the final transformer layer to generate a scalar reward estimate, following the setup in \cite{ziegler2020finetuning}. To minimize the variance of the learned reward, existing approaches center the reward predictions so that  $\mathbb{E}_{x,y\sim \mathcal{D}}\left[r_\phi(x, y)\right] = 0$ across all $x$.

\textbf{RL Fine-Tuning Phase}: In the subsequent reinforcement learning phase, the estimated reward function provides evaluative signals for model outputs. Building on previous work~\citep{jaques2017sequence, jaques2020human}, optimization is posed as

\begin{equation}\label{eq:RL}
\max_{\pi_{\theta}}  \mathbb{E}_{x\sim \mathcal{D}, y\sim \pi_{\theta}(y \mid x)}\bigl[r_{\phi}(x, y)\bigr] - \beta\mathbb{D}_{\textrm{KL}}\bigl[\pi_{\theta}(y\mid x)\mid \mid \pi_\text{ref}(y\mid x)\bigr],
\end{equation}

where the parameter $\beta$ regulates divergence from the reference policy $\pi_\text{ref}$, which is typically chosen as the original SFT policy $\pi^\text{SFT}$. In implementation, the initial weights of the language model policy $\pi_\theta$ match those of $\pi^\text{SFT}$. This regularization term serves to anchor the updated policy near the data regime where the reward model remains reliable, as well as promoting diversity in outputs and avoiding collapse onto a single output. As text generation is discrete, the resulting objective is non-differentiable and typically requires reinforcement learning for optimization. A common method \citep{ziegler2020finetuning, stiennon2022learning, bai2022training, ouyang2022training} is to define the reward as ${r(x, y) = r_{\phi}(x, y) -\beta (\log \pi_{\theta}(y\mid x) - \log \pi_\text{ref}(y\mid x))}$ and optimize using the PPO algorithm \cite{schulman2017proximal}. 

\section{Direct Preference Optimization}\label{sec:DPO}

Acknowledging the obstacles involved with scaling reinforcement learning techniques to the fine-tuning of large language models, we seek to introduce a streamlined method for policy optimization that directly utilizes preference data. Rather than following the standard RLHF paradigm of learning an explicit reward and performing RL-based optimization, we capitalize on a specific parameterization of the reward function that allows for direct, closed-form computation of the associated optimal policy, removing the necessity for a separate RL loop. As elaborated below, our main contribution lies in exploiting a closed-form relation between reward functions and their corresponding optimal policies, thereby enabling the conversion of an objective over rewards into one over policies. This reparameterization approach eliminates the need to train a separate reward predictor, yet remains consistent with widely-adopted preference modeling approaches like the Bradley-Terry framework. Effectively, the policy network unifies the roles of both language generation

\textbf{Deriving the DPO objective.} Our starting point is the standard RL objective previously explored in the literature, as presented in Eq.~\ref{eq:RL}, where we consider a generic reward function $r$. Building on established results~\citep{peters2007reinforcement, peng2019advantage, korbak2022reinforcement, go2023aligning}, the solution to the KL-regularized reward maximization objective in Eq.~\ref{eq:RL} is given by:

\begin{equation}\label{eq:op_policy}
    \pi_r(y\mid x) = \frac{1}{Z(x)}\pi_\text{ref}(y\mid x)\exp\left(\frac{1}{\beta}r(x, y)\right),
\end{equation}

where the normalization constant $Z(x)$ is defined as $Z(x) =\sum_{y}\pi_\text{ref}(y\mid x)\exp\left(\frac{1}{\beta}r(x, y)\right)$. A detailed derivation is provided in Appendix \ref{app:derivation1}. However, even when employing an MLE-trained reward estimator $r_\phi$ to approximate the true reward $r^*$, the computation of $Z(x)$ remains a bottleneck due to its intractability~\citep{korbak2022reinforcement, go2023aligning}. Despite this difficulty, we may reformulate Eq.~\ref{eq:op_policy} to characterize the reward using the optimal policy $\pi_r$, the reference policy $\pi_\text{ref}$, and the uncomputed $Z(\cdot)$. To do this, we first apply the logarithm to both sides of Eq.~\ref{eq:op_policy} and rearrange terms to obtain:

\begin{equation}\label{eq:main_eq}
    r(x,y) =\beta \log \frac{\pi_r(y\mid x)}{\pi_\text{ref}(y\mid x)} + \beta \log Z(x).
\end{equation}

This transformation can be readily applied to the true reward $r^*$ and the associated optimal policy $\pi^*$. Crucially, the Bradley-Terry framework depends solely on reward differences for two candidates, specifically, ${p^*(y_1 \succ y_2 \mid x) = \sigma(r^*(x, y_1) - r^*(x, y_2))}$. Replacing $r^*(x,y)$ in the preference likelihood with its equivalent in Eq.~\ref{eq:main_eq}, the normalization constants $Z(x)$ are eliminated, yielding a preference model that depends exclusively on $\pi^*$ and $\pi_\text{ref}$. Therefore, the optimal RLHF policy $\pi^*$ under the Bradley-Terry model satisfies:

\begin{equation}\label{eq:objective}
    p^*(y_1\succ y_2 \mid x)=\frac{1}{1 + \exp\left(\beta \log \frac{\pi^*(y_2\mid x)}{\pi_\text{ref}(y_2\mid x)} - \beta \log \frac{\pi^*(y_1\mid x)}{\pi_\text{ref}(y_1\mid x)}\right)}
\end{equation}

For the details of this derivation, see Appendix~\ref{app:derivation2}. Although Eq.~\ref{eq:objective} adopts the Bradley-Terry model, parallel derivations exist for the Plackett-Luce models~\citep{plackett1975analysis, luce2012individual} as described in Appendix~\ref{app:plackett_luce_models}.

Expressing human preference probabilities in terms of the optimal policy, instead of the reward, now enables us to propose a maximum likelihood training procedure for a parameterized policy $\pi_\theta$. Mirroring the maximum likelihood principle applied to reward models (see Eq.~\ref{eq:reward_model}), this yields the following policy learning objective:

\begin{equation}\label{eq:optimum_model}
    \mathcal{L}_\text{DPO}(\pi_{\theta}; \pi_\text{ref}) = -\mathbb{E}_{(x, y_w, y_l)\sim \mathcal{D}}\left[\log \sigma \left(\beta \log \frac{\pi_{\theta}(y_w\mid x)}{\pi_\text{ref}(y_w\mid x)} - \beta \log \frac{\pi_{\theta}(y_l\mid x)}{\pi_\text{ref}(y_l\mid x)}

In this approach, we estimate an implicit reward via a distinct parameterization, where the optimal policy corresponds directly to $\pi_\theta$. Furthermore, because our methodology mirrors the fitting of a reparameterized Bradley-Terry model, it inherits specific theoretical guarantees—such as consistency under appropriate conditions regarding the distribution of preference data \cite{bong2022generalized}. We explore these theoretical aspects of DPO and compare them to related methods in Section~\ref{sec:theory}.

\textbf{Mechanistic Interpretation of the DPO Update} To clarify how DPO operates, it is instructive to inspect the gradient of the loss function $\mathcal{L}_\text{DPO}$. The gradient with respect to $\theta$ is given by:

\begin{multline*}\label{eq:gradient}
    \nabla_\theta \mathcal{L}_\text{DPO}(\pi_\theta;\pi_\text{ref}) = \\ -\beta\mathbb{E}_{(x, y_w, y_l) \sim \mathcal{D}} \bigg[\underbrace{\sigma(\hat{r}_\theta(x, y_l) - \hat{r}_\theta (x, y_w))}_\text{greater emphasis when reward ranking is inaccurate}\bigg[\underbrace{\nabla_\theta\log \pi(y_w \mid x)}_\text{favor $y_w$} - \underbrace{\nabla_\theta\log\pi(y_l \mid x)}_\text{discourage $y_l$}\bigg]\bigg],
\end{multline*}

with $\hat{r}_\theta(x, y) = \beta \log \frac{\pi_\theta(y \mid x)}{\pi_\text{ref}(y \mid x)}$ serving as the latent reward formulated by the language model $\pi_\theta$ relative to a reference $\pi_\text{ref}$ (details in Section~\ref{sec:theory}). Conceptually, the loss gradient $\mathcal{L}_\text{DPO}$ nudges the model to assign higher probabilities to preferred completions $y_w$ and lower probabilities to less-preferred completions $y_l$. Crucially, the influence of each training instance is modulated by the extent to which the implicit reward model $\hat{r}_\theta$ overestimates the less-preferred option, weighted by $\beta$—effectively reflecting the magnitude of reward misordering and incorporating the KL constraint's impact. Our empirical results highlight the necessity of this weighting mechanism; omitting the coefficient leads to undesirable model collapse (see Appendix Table~\ref{tab:unlikelihood_generations}).

\textbf{DPO Overview.} The typical DPO workflow involves two primary steps: First, for each input prompt $x$, generate candidate completions $y_1, y_2 \sim \pi_\text{ref}(\cdot \mid x)$, and then collect human preference judgments to form an offline dataset $\mathcal{D} = \{x^{(i)}, y_w^{(i)}, y_l^{(i)}\}_{i=1}^N$. Second, optimize the language model parameters $\pi_\theta$ by minimizing the DPO loss $\mathcal{L}_\text{DPO}$ given the reference model $\pi_\text{ref}$, dataset $\mathcal{D}$, and a specified $\beta$ value.

In real-world settings, it is often preferable to leverage existing public preference datasets rather than produce new samples and acquire new preference labels. If such datasets have been generated using a supervised fine-tuned model $\pi^\text{SFT}$, we initialize $\pi_\text{ref} = \pi^\text{SFT}$. When no such model $\pi^\text{SFT}$ exists, $\pi_\text{ref}$ is initialized instead by maximizing the likelihood of the preferred outputs $(x, y_w)$, i.e., $\pi_\text{ref} = \argmax_{\pi}\mathbb{E}_{x, y_w \sim \mathcal{D}}\left[\log \pi(y_w \mid x)\right]$. This initialization addresses the possible mismatch between the unobservable true reference distribution and the operational $\pi_\text{ref}$ within DPO. Additional implementation specifics and information on hyperparameters can be found in Appendix~\ref{app:implementation}.

\section{Theoretical Analysis of DPO} This section provides a deeper theoretical grounding for DPO, offering interpretive insights and comparing its strengths to traditional actor-critic RLHF techniques such as PPO~\cite{schulman2017proximal}.

\label{sec:theory}

\subsection{Your Language Model Is Secretly a Reward Model} DPO effectively combines reward modeling and policy learning within a single maximum likelihood framework, sidestepping the explicit construction of a reward model and separate reinforcement learning procedures. The optimization problem in Eq. \ref{eq:main_eq} corresponds to the Bradley-Terry model, where rewards are parameterized as $r^*(x, y) = \beta \log\frac{\pi^*_\theta(y \mid x)}{\pi_\text{ref}(y \mid x)}$. Optimizing our parameterized policy $\pi_{\theta}$ under this view is equivalent to optimizing the reward model as described by Eq. \ref{eq:reward_model}, after applying a suitable change of variables. In what follows, we formalize the theoretical justification for this reparameterization, demonstrate that it imposes no significant limitations on the range of reward models that can be represented, and establish that it supports exact recovery of the optimal policy. We start by introducing a formal equivalence relation over reward functions.

\begin{definition}
Two reward functions $r(x, y)$ and $r'(x, y)$ are said to be equivalent if and only if ${r(x, y)-r'(x, y) = f(x)}$ for some function $f$.
\end{definition}

It is straightforward to verify that this constitutes an equivalence relation, thereby dividing the set of reward functions into distinct equivalence classes. This observation leads directly to two key lemmas:

\begin{lemma}\label{lemma:same_prefrence}
Within the Plackett-Luce family—specifically, the Bradley-Terry model—two reward functions belonging to the same equivalence class give rise to identical preference distributions.
\end{lemma}

\begin{lemma}\label{lemma:same_policy}
    Any two reward functions from the same equivalence class result in the same optimal policy when considering the constrained RL setting.
\end{lemma}

The justifications for these statements are elementary and can be found in Appendix \ref{app:lemma1}. The first lemma highlights a well-documented issue of under-specification present in models from the Plackett-Luce family \cite{plackett1975analysis}. This ambiguity requires the introduction of additional identifiability constraints in order to guarantee meaningful MLE estimation according to Eq. \ref{eq:reward_model} \cite{bong2022generalized}. The second lemma clarifies that the specific choice of representative reward function within an equivalence class does not affect the resulting optimal policy, thus for our primary aim, it suffices to recover any member from the correct optimal equivalence class. The following Theorem, whose proof appears in Appendix~\ref{app:thm1}, establishes a useful representation:

\begin{theorem}\label{thm:main}
    Assuming mild technical conditions, every reward equivalence class compatible with Plackett-Luce (and specifically Bradley-Terry) models admits the reparameterization ${r(x, y) = \beta \log \frac{\pi(y\mid x)}{\pi_\text{ref}(y\mid x)}}$ for some model $\pi(y\mid x)$, given a reference model $\pi_\text{ref}(y \mid x)$.
\end{theorem}

\begin{sproof}
    Take any reward function $r(x, y)$, which determines its corresponding optimal policy $\pi_r(y \mid x)$ by Eq. \ref{eq:op_policy}. We now demonstrate that within the equivalence class of $r$, one can always find a function that takes the above reparameterized form. Define the following projection:
\begin{equation}
    f(r; \pi_\text{ref}, \beta)(x, y) = r(x, y) - \beta\log\sum_{y}\pi_\text{ref}(y\mid x)\exp\left(\frac{1}{\beta}r(x, y)\right)
\end{equation}
Here, $f$ serves to normalize $r$ using the log-partition function over $\pi_r$. The normalization term depends solely on $x$, thus $f(r; \pi_\text{ref}, \beta)(x, y)$ stays within the same equivalence class as $r(x, y)$. Substituting $r$ by the right-hand side of Eq.~\ref{eq:main_eq} (valid for any $r$), yields $f(r; \pi_\text{ref}, \beta)(x, y) = \beta \log \frac{\pi_r(y\mid x)}{\pi_\text{ref}(y\mid x)}$. Thus, $f$ constructs a function in the same equivalence class with the desired parameterization, ensuring no loss of generality from adopting this specific reparameterized reward model.
\end{sproof}

An alternative perspective on Theorem~\ref{thm:main} is that it precisely determines the specific reward function within each equivalence class that is chosen by the DPO reparameterization; namely, the reward function that fulfills:

\begin{equation}\label{eq:lag_p}
     \sum_{y}\underbrace{\pi_\text{ref}(y\mid x)\exp\left(\frac{1}{\beta}r(x, y)\right)}_{=\pi(y\mid x)\text{, using Thm.~\ref{thm:main} reparam.}} = 1,
\end{equation}

In other words, $\pi(y\mid x)$ constitutes a proper probability distribution, with non-negative values that sum to one. Building on Eq.~\ref{eq:op_policy}, we observe that Eq.~\ref{eq:lag_p} represents the partition function of the optimal policy that emerges from the reward function $r(x, y)$. The central insight underpinning the DPO algorithm is that by introducing constraints to the underdetermined Plackett-Luce (and specifically Bradley-Terry) family of preference models, one can retain the full expressive power of representable reward functions, while rendering the optimal policy from Eq.~\ref{eq:op_policy} analytically manageable for any prompt $x$.

\subsection{Instability of Actor-Critic Algorithms} Our proposed framework further allows for the analysis of instabilities often encountered in standard actor-critic approaches for RLHF, such as PPO. We adopt the RLHF protocol and concentrate on the RL fine-tuning phase as detailed in Section \ref{section:prelims}. The structure of our approach aligns with the control as inference framework \cite{levine2018reinforcement}, as applied to the constrained RL setup introduced in \ref{eq:RL}. Assuming a parameterized policy $\pi_{\theta}(y\mid x)$, our aim is to minimize the KL divergence $\mathbb{D}_{\text{KL}}[\pi_{\theta}(y|x) \mid \mid \pi^*(y\mid x)]$, where $\pi^*$ denotes the optimal policy defined in Eq. \ref{eq:optimum_model} and derived from reward $r_{\phi}(y, x)$. Through a straightforward derivation, this yields the following optimization objective:

\begin{equation}\label{eq:AC}
    \max_{\pi_{\theta}}\mathbb{E}_{\pi_{\theta}(y\mid x)}\bigg[\underbrace{r_{\phi}(x, y) -\beta\log\sum_{y}\pi_\text{ref}(y\mid x)\exp\left(\frac{1}{\beta}r_{\phi}(x, y)\right)}_{f(r_{\phi}, \pi_\text{ref}, \beta)} - \underbrace{\beta\log\frac{\pi_{\theta}(y\mid x)}{\pi_\text{ref}(y\mid x)}}_{\text{KL}}\bigg]
\end{equation}

This objective is identical to that optimized in previous studies 
\citep{ziegler2020finetuning, stiennon2022learning, bai2022training, ouyang2022training}, which employ a DPO-equivalent reward for the reward family $r_{\phi}$. Within this context, the normalization component in $f(r_{\phi}, \pi_\text{ref}, \beta)$ can be viewed as the soft value function associated with the reference policy $\pi_\text{ref}$. Although this term has no influence on the final optimal policy, omitting it can result in elevated variance in the policy gradient, thereby destabilizing the learning process. This normalization can be incorporated via a learned value function, though this approach itself may introduce optimization challenges. Another strategy employed in prior literature is to normalize the reward using a baseline corresponding to human completion, which acts as a single-sample Monte Carlo approximation of the normalization. Distinctively, the DPO reparameterization yields a reward function inherently free from the necessity for such baselines.

\section{Experiments}
We proceed by empirically assessing DPO’s capacity to train policies from preference data. Our initial investigation considers a controlled environment for text generation, in which we examine how effectively DPO navigates the trade-off between reward maximization and minimization of KL-divergence relative to the reference policy, particularly when compared to established preference learning methods like PPO. Subsequently, we test DPO on larger-scale models and on challenging RLHF tasks, such as summarization and conversational agents. Our observations indicate that, with minimal tuning of hyperparameters, DPO generally matches or exceeds the performance of robust baselines like RLHF with PPO, as well as the performance of strategies that return the best out of $N$ samples according to a learned reward model. Prior to discussing these empirical findings, we outline the experimental procedures used; further methodological information can be found in Appendix~\ref{app:exp_details}.

\textbf{Tasks.} Our study investigates three distinct open-ended text generation challenges. For every task, the learning algorithms derive a policy based on a dataset of preferences denoted as $\mathcal{D}=\bigl\{x^{(i)}, y_w^{(i)}, y_l^{(i)}\bigr\}_{i=1}^N$. The first task, \textbf{controlled sentiment generation}, uses $x$ as a prompt drawn from an IMDb movie review prefix \cite{maas-EtAl:2011:ACL-HLT2011}. Here, the aim is for the policy to generate $y$ that conveys a positive sentiment. To enable rigorous evaluation, we synthetically create preference pairs for model outputs by employing a pre-trained sentiment classifier, enforcing $p(\text{positive}\mid x,y_w)>p(\text{positive}\mid x,y_l)$. For SFT, GPT-2-large is fine-tuned until it converges using the training portion of the IMDB dataset; additional training specifics can be found in App~\ref{app:sentiment_details}. The second task, \textbf{summarization}, considers $x$ as a Reddit forum post, with the model required to generate $y$, a summary capturing the main points of the post. Building on earlier approaches, the Reddit TL;DR summarization dataset \citep{volske-etal-2017-tl} is employed, supplemented with human preference data curated by \citeauthor{stiennon2022learning}. The SFT model is trained on summaries written by humans for forum posts,\footnote{\url{https://huggingface.co/CarperAI/openai_summarize_tldr_sft}} using the TRLX \citep{leandro_von_werra_2023_7790115} framework for RLHF. The associated preference labels were obtained by \citeauthor{stiennon2022learning} from samples produced by a separate, but similarly fine-tuned, SFT model. The third task, \textbf{single-turn dialogue}, designates $x$ as a user's question or statement, ranging from technical scientific inquiries to personal advice requests. The objective is for the policy to produce an informative and engaging reply $y$. This experiment draws from the Anthropic Helpful and Harmless dialogue dataset \citep{bai2022training}, which comprises 170k conversational exchanges between humans and an AI assistant. Each conversation is concluded by two responses from a large (though unspecified) language model, accompanied by a human-annotated label indicating the preferred answer. Since no dedicated SFT model exists for this domain, we fine-tune a pre-trained language model using only those completions judged as preferred to create the SFT model for dialogue.

\textbf{Evaluation.} Our study employs two primary evaluation methodologies. For controlled sentiment generation, where the true reward function (a sentiment classifier) is accessible, we examine each algorithm’s capability to optimize the constrained reward maximization objective by charting the frontier between attained reward and KL-divergence from the reference policy; this allows for precise quantification. Conversely, since the reward function is unavailable in real-world scenarios, we resort to measuring the algorithms' \textit{win rate} when compared to a baseline policy. In these cases, GPT-4 serves as a stand-in for human judgment to assess the quality of summaries and the helpfulness of responses in summarization and single-turn dialogue tasks, respectively. For summarization, the baseline is the reference summary from the test set, while for dialogue, the baseline is the preferred response as indicated in the test data. While prior research suggests large language models (LMs) can surpass traditional automatic metrics as evaluators \citep{Chen2023ExploringTU}, we additionally perform a human evaluation to support our reliance on GPT-4 as a judge, detailed in Sec.~\ref{sec:human-judgments}. Our findings indicate a high degree of correlation between GPT-4 and human ratings; notably, agreement between humans and GPT-4 matches or surpasses inter-human annotator agreement.

\textbf{Methods.} Alongside DPO, we benchmark a range of existing strategies for training language models aligned with human preferences. For the summarization task, we examine zero-shot prompting with \textbf{GPT-J} \citep{gpt-j}, and in the dialogue setting, we utilize 2-shot prompting with \textbf{Pythia-2.8B} \citep{biderman2023pythia}. Furthermore, we include the \textbf{SFT} model and \textbf{Preferred-FT}, which involves supervised fine-tuning on the chosen response $y_w$—originating from SFT in sentiment control and summarization, or from a generic LM in dialogue. The \textbf{Unlikelihood} approach \citep{welleck2019neural} represents another pseudo-supervised method; it updates the policy by maximizing likelihood of $y_w$ and explicitly penalizing $y_l$, modulated by a tunable coefficient $\alpha\in[0,1]$ on the unlikelihood term. We additionally evaluate \textbf{PPO} \citep{schulman2017proximal} utilizing a reward model derived from preference annotations, and \textbf{PPO-GT}, an oracle version trained directly on the true reward function present in the controlled sentiment scenario. Within the sentiment experiments, we test both a standard PPO-GT implementation \cite{leandro_von_werra_2023_7790115} and an enhanced version that normalizes rewards and further tunes hyperparameters to maximize performance; similar adjustments are applied when running conventional PPO with learned reward functions. Lastly, we implement the \textbf{Best of $N$} strategy: $N$ responses are sampled from the SFT model (or Preferred-FT in the dialogue setting), and the response with the highest predicted reward, based on a reward model trained on preference data, is selected. Despite strong empirical performance, this approach decouples reward modeling from PPO optimization but is rendered inefficient for even modest $N$, given the necessity to generate $N$ completions per test instance.

\subsection{How well can DPO optimize the RLHF objective?}

In standard RLHF methodologies, the primary goal is to maximize reward while simultaneously constraining the policy’s divergence from a reference distribution via a KL penalty. Accordingly, when evaluating different approaches, it is crucial to assess both the achieved reward and the KL distance; a modest increase in reward at the expense of a large increase in KL may not be favorable. Figure~\ref{fig:frontier-tldr-main} illustrates the tradeoff frontier between reward and KL across several methods within the sentiment domain. For each method, we conduct a set of training runs, varying the hyperparameters governing policy conservatism: target KL $\in\{3,6,9,12\}$ for PPO, $\beta \in \{0.05,0.1,1,5\}$, and $\alpha\in\{0.05,0.1,0.5,1\}$ for unlikelihood, while using distinct random seeds for preferred-FT. In total, this comprehensive sweep covers 22 separate runs. Every 100 training steps up to convergence, we evaluate the models by calculating the mean reward—according to the actual reward function—and the average sequence-level KL\footnote{This refers to the aggregate of KL divergences computed at each timestep.} compared to the reference policy, i.e., $\text{KL}\left(\pi\mid \mid \pi_\text{ref}\right)$. Our results demonstrate that DPO establishes the most favorable frontier, delivering the highest rewards with relatively low KL. This finding stands out for several reasons. Notably, even though DPO and PPO share the underlying objective, DPO is far more efficient—its reward/KL performance consistently exceeds PPO's. Moreover, DPO yields a superior reward-KL frontier in comparison to PPO, \emph{even under conditions where PPO has access to ground-truth rewards} (PPO-GT).

\subsection{Can DPO scale to real preference datasets?}
\label{sec:dpo-real-datasets}
We proceed to investigate how well DPO performs when fine-tuned on real-world preference datasets, focusing on tasks such as summarization and single-turn dialogue. In the context of summarization, automatic measures like ROUGE often show weak alignment with human judgment~\citep{stiennon2022learning}, and previous studies have demonstrated that employing PPO to fine-tune language models on human preferences results in more compelling summaries. For our comparison, we generate model outputs using the test partition of the TL;DR summarization dataset and assess performance via the average win rate when pitted against reference summaries from the test data. We sample model completions across a range of temperatures (from 0.0 to 1.0), and present the corresponding win rates in Figure~\ref{fig:frontier-tldr-main} (right). DPO, PPO, and Preferred-FT all leverage the identical GPT-J SFT model as a starting point\footnote{\url{https://huggingface.co/CarperAI/openai_summarize_tldr_sft}}. Our experiments reveal that DPO achieves a win rate close to 61\% at a temperature setting of 0.0, surpassing PPO, which records about 57\% at its optimal temperature of 0.0. Additionally, DPO attains a higher peak win rate than the best of $N$ baseline. It should be highlighted that no significant hyperparameter tuning was performed for DPO’s $\beta$, suggesting that these outcomes may not fully capture DPO's upper-bound capabilities. Furthermore, DPO exhibits considerably greater resilience to changes in sampling temperature compared to PPO, whose performance can decline to that of the base GPT-J model at higher temperatures. Notably, Preferred-FT does not show marked gains over the initial SFT model. Direct comparisons between DPO and PPO in human preference experiments (see Section~\ref{sec:human-judgments}) indicate that DPO samples, generated at a temperature of 0.25, are chosen 58\% of the time over PPO samples produced at the same temperature.

For single-turn dialogue, we assess the various approaches using the subset of the Anthropic HH dataset test split \citep{bai2022training} that consists of one human-assistant exchange. In our GPT-4 evaluation, we take the completions marked as preferred in the test set as ground truth references, calculating win rates for each method accordingly. Since there is not an established SFT model available for this task, we initiate our experiments with the Pythia-2.8B pre-trained model, apply Preferred-FT to fit a reference model to the selected completions so its outputs align with the relevant data distribution, and subsequently train with DPO. For benchmarking, we include both the strongest outcome among 128 Preferred-FT completions (having observed that the Best of $N$ metric saturates at $N = 128$ for this evaluation—see Appendix Figure~\ref{fig:best-of-n}) and a 2-shot prompted Pythia-2.8B base model. Across methods, DPO matches or outperforms the others at their optimal sampling temperatures. We additionally evaluate an RLHF model trained with PPO on the Anthropic HH dataset \footnote{\url{https://huggingface.co/reciprocate/ppo_hh_pythia-6B}} from a recognized repository \footnote{\url{https://github.com/CarperAI/trlx/tree/main/examples/hh}}, but are unable to identify any prompt or temperature that yields results surpassing the base Pythia-2.8B’s performance. Given our observations with TL;DR and the alignment of optimization objectives, we treat Best of 128 as an approximate indicator for PPO performance. Overall, DPO emerges as the only method that both achieves greater efficiency and surpasses the preferred completions in the Anthropic HH set, while also matching or exceeding the results of the more computationally intensive Best of 128 baseline. Lastly, as illustrated in Figure~\ref{fig:dialogue-main}, DPO attains its peak performance rapidly during training.

\subsection{Generalization to a new input distribution}

To more thoroughly assess how PPO and DPO perform when encountering distributional changes, we apply the policies derived from our Reddit TL;DR summarization study to an out-of-distribution dataset—specifically, news articles from the CNN/DailyMail test set \citep{nallapati-etal-2016-abstractive}. We use the top-performing sampling temperatures identified from the TL;DR experiments (0 and 0.25). The corresponding results are displayed in Table~\ref{tab:ood}. The GPT-4 win rate relative to ground-truth summaries was determined using the same GPT-4 (C) prompt as in the Reddit TL;DR evaluation, except the term ``forum post'' was substituted with ``news article''. Notably, DPO continues to show a marked advantage over PPO on this alternate dataset. These findings suggest that DPO policies demonstrate a generalization capability on par with, if not exceeding, that of PPO policies, even without leveraging the supplementary unlabeled Reddit TL;DR prompts that PPO employs.

\subsection{Validating GPT-4 judgments with human judgments}
\label{sec:human-judgments}
We undertake a human evaluation study to examine the consistency between GPT-4’s assessments and those of human annotators, focusing on outputs from the TL;DR summarization task and utilizing two distinct GPT-4 prompts. The \textbf{GPT-4 (S)} (simple) prompt inquires which summary better encapsulates the essential content of the original post. The \textbf{GPT-4 (C)} (concise) prompt extends this by also querying which summary is more concise; this is important because GPT-4, when following the \textbf{GPT-4 (S)} prompt, shows a tendency to prefer lengthier and more repetitive outputs than human evaluators do. The full prompt texts can be found in Appendix~\ref{app:prompts}. For our comparison, we select three systems representing a spectrum of output quality: the strongest (DPO, temperature 0.25), the weakest (PPO, temperature 1.0), and an intermediate model (SFT, temperature 0.25), each benchmarked against PPO with greedy decoding (its optimal temperature). Our results indicate that, regardless of prompt, GPT-4 aligns with human preferences at a rate comparable to the inter-annotator agreement among humans themselves. This suggests that GPT-4 serves as a reliable surrogate for human evaluations, although our collection of multiple human judgments was limited to the DPO and PPO-1 cases due to rater constraints. On the whole, the \textbf{GPT-4 (C)} prompt yields win rates that more closely mirror human preferences, which is why we use it for the primary analyses presented in Section~\ref{sec:dpo-real-datasets}. More detailed information on the human study—covering the web interface and the roster of human raters—is available in Appendix~\ref{app:human-study}.

\section{Discussion}
Preference-based learning offers a robust and scalable strategy for developing language models that are both competent and aligned with human values. In this work, we have presented DPO, a straightforward approach that enables training language models from preferences without the need for reinforcement learning. Instead of recasting the preference learning challenge within a conventional RL framework to leverage standard RL techniques, DPO establishes a correspondence between language model policies and reward functions, thereby allowing direct optimization towards human preferences using a cross-entropy loss function. This is achieved without engaging reinforcement learning or sacrificing generality. Remarkably, DPO matches or exceeds the performance of existing RLHF methods such as those based on PPO, often requiring minimal hyperparameter adjustment; thus, it significantly lowers the difficulty associated with preference-based language model training.

\textbf{Limitations \& Future Work.} Several avenues for future investigation emerge from our findings. One open question concerns the out-of-distribution generalization capacity of the DPO policy, particularly in comparison to approaches that rely on explicit reward functions; while preliminary evidence points to DPO generalizing on par with PPO-driven methods, broader empirical analyses are necessary. Additionally, it is worth exploring whether self-labeling mechanisms utilizing the DPO policy could leverage unlabeled data as effectively. Another pertinent issue is understanding how phenomena such as reward over-optimization appear within the DPO framework, and whether the observed modest decline in performance in Figure~\ref{fig:dialogue-main}-right reflects this effect. Furthermore, our present work is limited to models of up to 6B parameters, and scaling DPO to handle significantly larger, cutting-edge architectures remains a promising topic for future study. On the evaluation front, we observed that GPT-4-derived win rates are sensitive to the prompt formulation; subsequent research may delve into optimal methods for procuring reliable assessments from automated evaluators. Ultimately, the potential applications for DPO span well beyond language modeling, offering promising prospects for training generative systems across diverse modalities.